\seccionanexo{Anexo B: Validación y Confiabilidad de los Instrumentos}

La validación de los instrumentos (cuestionarios cerrados con escala Likert de los Anexos A y C) se realizó mediante juicio de experto y prueba piloto, sin cálculo de alfa de Cronbach por alcance descriptivo y muestreo por conveniencia:

\begin{itemize}
    \item \textbf{Juicio de experto (Tutor Guía):} El Ing. Juan Vladimir Ramirez Flores evaluó claridad, pertinencia y coherencia de los ítems respecto a las 3 brechas (recovery, documental, paralelos) y a las 7 variables del estudio. Criterios: escala 1--4 (1 no pertinente, 4 muy pertinente). Observaciones: simplificar redacción ítems 3--4 (evitar tecnicismos “trazabilidad”) y añadir instrucción de anonimato (C.I. como ****). Fecha de validación: febrero 2026. Dictamen: \textit{Aplicable con ajustes menores} (promedio 3.7/4).
    \item \textbf{Prueba piloto:} Aplicada a 5 estudiantes postulantes no incluidos en la muestra final (2 de Sistemas Informáticos, 1 de Electrónica, 1 de Gastronomía, 1 de Construcción Civil), con tiempo promedio de respuesta de 4 min y sin incidencias de comprensión. Ajuste resultante: ítem “entrega y verificación … en secretaría” → “con los tutores de curso” para reflejar el flujo V2.0.
    \item \textbf{Confiabilidad operativa:} No se aplicó prueba inferencial; la contrastación es descriptiva y de auditoría transaccional (ver Anexo D). Se estableció criterio de validez: latencia menor a 1.2 s en al menos 95\% de peticiones, 0 sobrecupos con guardado seguro, papeles completos 100\%, facilidad de uso de al menos 90\% y errores menores a 5\% (ver Cap.~5).
\end{itemize}

\begin{table}[htbp]
    \centering
    \caption{Resultados de la prueba piloto (n=5, anonymized).}
    \label{tab:piloto_validacion}
    \begin{tabularx}{\textwidth}{C C C}
        \toprule
        \textbf{Dimensión} & \textbf{Comprensión (\%)} & \textbf{Tiempo medio} \\
        \midrule
        Usabilidad y Diseño Móvil & 100 & 1.1 min \\
        Autoservicio y Recovery & 80 & 1.3 min \\
        Matriculación Documental & 100 & 0.8 min \\
        Organización Paralelos & 100 & 0.8 min \\
        \midrule
        \textbf{Total} & \textbf{95} & \textbf{4.0 min} \\
        \bottomrule
    \end{tabularx}
    \notatabla{Prueba piloto febrero 2026; datos anonimizados. Ajustes incorporados en Anexo C.}
\end{table}
